\section{Tabella dei Casi d'Uso}

\begin{longtable}{|p{3cm}|p{11cm}|}
\hline
\textbf{Attore} & \textbf{Casi d'uso} \\
\hline
\endfirsthead
\hline
\textbf{Attore} & \textbf{Casi d'uso} \\
\hline
\endhead
\hline
\endfoot
\hline
\endlastfoot

Utente non registrato &
\begin{itemize}[leftmargin=*,nosep,nolistsep]
    \item Registrazione al sistema tramite email
    \item Visualizzazione informazioni pubbliche dell'applicazione
    \item Accesso alla pagina di login
\end{itemize} \\
\hline

Utente registrato (USER) &
\begin{itemize}[leftmargin=*,nosep,nolistsep]
    \item Login/logout al sistema
    \item Visualizzazione lista bug con filtri (tipo, stato, priorità, label)
    \item Visualizzazione dettagli singolo bug
    \item Creazione nuovo bug
    \item Modifica bug propri o assegnati
    \item Aggiunta commenti ai bug
    \item Upload allegati immagini ai bug
    \item Assegnazione bug ad altri utenti
    \item Archiviazione bug risolti
    \item Marcatura bug come duplicati
    \item Visualizzazione storico modifiche bug
    \item Gestione profilo personale
    \item Cambio password
\end{itemize} \\
\hline

Utente ReadOnly &
\begin{itemize}[leftmargin=*,nosep,nolistsep]
    \item Login al sistema
    \item Visualizzazione lista bug con filtri
    \item Visualizzazione dettagli singolo bug
    \item Visualizzazione commenti e allegati
    \item Visualizzazione storico modifiche
    \item Visualizzazione profilo proprio (sola lettura)
\end{itemize} \\
\hline

Admin &
\begin{itemize}[leftmargin=*,nosep,nolistsep]
    \item Tutti i casi d'uso dell'Utente registrato
    \item Gestione completa utenti (visualizzazione, modifica ruoli, disabilitazione)
    \item Visualizzazione metriche di sistema
    \item Generazione report mensili
    \item Esportazione dati in formato Excel
    \item Accesso completo a tutti i bug del sistema
    \item Gestione label e categorie
    \item Configurazione parametri di sistema
\end{itemize} \\
\hline

Sistema &
\begin{itemize}[leftmargin=*,nosep,nolistsep]
    \item Invio notifiche automatiche su cambiamenti bug
    \item Archiviazione automatica bug dopo periodo di inattività
    \item Generazione automatica storico modifiche
    \item Validazione automatica allegati (tipo e dimensione)
    \item Pulizia automatica dati obsoleti
    \item Backup automatico dati
\end{itemize} \\
\hline
\caption{Tabella completa dei casi d'uso per BugBoard26}
\label{tab:casi_uso}
\end{longtable}

\subsection{Descrizione dettagliata dei casi d'uso}

\subsubsection{Utente non registrato}
\begin{itemize}
    \item \textbf{Registrazione al sistema tramite email}: L'utente può creare un nuovo account fornendo email, password e informazioni personali. Il sistema invia una email di conferma per verificare l'indirizzo.
    \item \textbf{Visualizzazione informazioni pubbliche}: Accesso alla pagina principale dell'applicazione con informazioni generali sul servizio di bug tracking.
    \item \textbf{Accesso alla pagina di login}: Possibilità di accedere al sistema utilizzando credenziali esistenti.
\end{itemize}

\subsubsection{Utente registrato (USER)}
\begin{itemize}
    \item \textbf{Login/logout al sistema}: Autenticazione tramite credenziali email/password con generazione di token JWT per sessioni sicure.
    \item \textbf{Visualizzazione lista bug con filtri}: Consultazione dell'elenco completo dei bug con possibilità di filtrare per tipo, stato, priorità, label e ricerca testuale.
    \item \textbf{Visualizzazione dettagli singolo bug}: Accesso completo alle informazioni di un bug specifico, inclusi descrizione, commenti, allegati e storico modifiche.
    \item \textbf{Creazione nuovo bug}: Inserimento di una nuova segnalazione con descrizione dettagliata, assegnazione priorità, tipo e label appropriate.
    \item \textbf{Modifica bug propri o assegnati}: Possibilità di aggiornare informazioni, stato e dettagli dei bug creati dall'utente o assegnati al suo account.
    \item \textbf{Aggiunta commenti ai bug}: Inserimento di commenti testuali per fornire aggiornamenti, soluzioni o richiedere chiarimenti.
    \item \textbf{Upload allegati immagini ai bug}: Caricamento di immagini (JPEG, PNG, GIF, WebP) fino a 5MB per documentare visivamente il problema.
    \item \textbf{Assegnazione bug ad altri utenti}: Trasferimento della responsabilità di un bug a un altro utente registrato nel sistema.
    \item \textbf{Archiviazione bug risolti}: Spostamento dei bug risolti in stato archiviato per mantenere pulita la lista attiva.
    \item \textbf{Marcatura bug come duplicati}: Collegamento di un bug a un altro esistente quando viene identificato come duplicato.
    \item \textbf{Visualizzazione storico modifiche bug}: Consultazione cronologica di tutte le modifiche apportate a un bug nel tempo.
    \item \textbf{Gestione profilo personale}: Modifica delle informazioni personali, preferenze e impostazioni dell'account.
    \item \textbf{Cambio password}: Aggiornamento della password di accesso con validazione di sicurezza.
\end{itemize}

\subsubsection{Utente ReadOnly}
\begin{itemize}
    \item \textbf{Login al sistema}: Accesso con credenziali per utenti con privilegi di sola lettura.
    \item \textbf{Visualizzazione lista bug con filtri}: Consultazione dell'elenco bug con possibilità di applicare filtri di ricerca.
    \item \textbf{Visualizzazione dettagli singolo bug}: Accesso in sola lettura alle informazioni complete di un bug.
    \item \textbf{Visualizzazione commenti e allegati}: Lettura di tutti i commenti e download/visualizzazione degli allegati associati.
    \item \textbf{Visualizzazione storico modifiche}: Consultazione della cronologia delle modifiche senza possibilità di intervento.
    \item \textbf{Visualizzazione profilo proprio}: Accesso limitato alle informazioni del proprio profilo utente.
\end{itemize}

\subsubsection{Admin}
\begin{itemize}
    \item \textbf{Tutti i casi d'uso dell'Utente registrato}: L'amministratore eredita tutte le funzionalità dell'utente standard.
    \item \textbf{Gestione completa utenti}: Creazione, modifica, disabilitazione e gestione ruoli degli utenti del sistema.
    \item \textbf{Visualizzazione metriche di sistema}: Accesso a dashboard con statistiche, grafici e KPI del sistema.
    \item \textbf{Generazione report mensili}: Creazione di report dettagliati sull'attività del sistema per periodo specificato.
    \item \textbf{Esportazione dati in formato Excel}: Estrazione completa dei dati in formato spreadsheet per analisi esterne.
    \item \textbf{Accesso completo a tutti i bug}: Possibilità di visualizzare e modificare qualsiasi bug del sistema indipendentemente dall'assegnazione.
    \item \textbf{Gestione label e categorie}: Creazione, modifica ed eliminazione delle etichette e categorie di classificazione.
    \item \textbf{Configurazione parametri di sistema}: Modifica delle impostazioni globali, limiti e comportamenti dell'applicazione.
\end{itemize}

\subsubsection{Sistema}
\begin{itemize}
    \item \textbf{Invio notifiche automatiche}: Meccanismo automatico di notifica via email o push per aggiornamenti sui bug seguiti.
    \item \textbf{Archiviazione automatica}: Processo schedulato per archiviare automaticamente bug inattivi dopo un periodo configurabile.
    \item \textbf{Generazione automatica storico}: Registrazione automatica di ogni modifica apportata ai bug nel database.
    \item \textbf{Validazione automatica allegati}: Controllo automatico di tipo, dimensione e sicurezza dei file caricati.
    \item \textbf{Pulizia automatica dati obsoleti}: Eliminazione periodica di dati temporanei, cache e informazioni non più necessarie.
    \item \textbf{Backup automatico dati}: Salvataggio automatico periodico del database e dei file allegati per disaster recovery.
\end{itemize}

